\documentclass[11pt,a4paper]{article}
\usepackage[utf8]{inputenc}
\usepackage[T1]{fontenc}
\usepackage[margin=1in]{geometry}
\usepackage{xcolor}
\usepackage{hyperref}
\usepackage{lmodern}
\definecolor{primarycolor}{RGB}{33,37,41}
\definecolor{secondarycolor}{RGB}{0,102,204}
\hypersetup{colorlinks=true,linkcolor=secondarycolor,urlcolor=secondarycolor,citecolor=secondarycolor}
\pagestyle{empty}
\setlength{\parindent}{0pt}
\setlength{\parskip}{0.85em}
\begin{document}
{\Large\bfseries\color{primarycolor} Marta Kowal-Brandt}\\[0.3em]
{\small \href{mailto:marta.kowalbrandt@example.com}{marta.kowalbrandt@example.com} | +48 555 303 303 | \href{https://gitlab.com/mkb-example}{gitlab.com/mkb-example}}

Corvale Mobility\\
\textbf{Subject: Cloud Infrastructure / DevOps Contractor -- Fleet Connectivity}

Dear Corvale Mobility Hiring Team,

I am applying for the Cloud Infrastructure / DevOps Contractor role in Fleet Connectivity. My route into DevOps is through hands-on ownership of CI/CD and test infrastructure: at Odrava Health Software, I own GitLab CI pipelines for three product teams supporting approximately 25 developers. Corvale's requirement to help backend developers deploy services that are easy to monitor and run connects directly with this work.

I reduced average pipeline time from 38 to 14 minutes and monthly CI runner costs by approximately 30\% through parallelisation, caching and changes to end-to-end execution. I also built merge-request environments that deploy the application, PostgreSQL and seeded data through Helm, then clean themselves up automatically. These environments let QA and product review features before merge. Alongside the implementation, I wrote the CI guide used in onboarding: making infrastructure understandable to colleagues is part of how I support the teams using it.

My personal infrastructure work extends that foundation into infrastructure as code, GitOps and observability. I provisioned a Hetzner VPS with Terraform and built a three-node k3s homelab using Argo CD, Prometheus, Grafana, Loki and cert-manager. My commercial Kubernetes work is at namespace level rather than cluster administration, and my Terraform and AWS experience comes from personal projects and learning, not production cloud operations. This is a transition into DevOps, grounded in delivery tooling and test environments rather than an established cloud architecture background.

For Corvale's vehicle-to-cloud platform, I would bring practical experience in repeatable deployments, diagnostic tooling and clear technical documentation. I would welcome a discussion about how those strengths could contribute to the Fleet Connectivity team as it takes the new infrastructure toward production. Thank you for considering my application.

Kind regards,\\
Marta Kowal-Brandt
\end{document}